\documentclass[fontset=fandol]{ctexart}
\usepackage[margin=16mm]{geometry}
\usepackage{gradbars}
\setlength{\parindent}{0pt}
\usegradbarslibrary{distribution}
\gradbarsdeclarelibrary{sample-columns}{table}{docs/libraries/gradbars-contrib-columns.sty}
\usegradbarslibrary{sample-columns}
\begin{document}
\section*{显式列适配与原始样本分布}
虚构复核数据；扩展库仅使用公开 API。
\par
\begin{gradtable}[style=adapter-report,width=\linewidth,label width=0pt]
 \gradcolumn[width=28mm]{project}{项目}
 \gradcolumn[type=review-score,mark=both,aggregate=mean]{score}{评分}
 \gradcolumn[type=review-stamp,width=24mm]{stamp}{固定盒}
 \gradrow{Atlas,80,80}\gradrow{Birch,0,0}\gradrow{Cedar,40,40}
 \gradsummary{均值}
\end{gradtable}\par\bigskip
\par
\begin{gradtable}[style=rug-report,width=\linewidth]
 \gradcolumn[width=28mm]{batch}{批次}
 \gradcolumn[type=boxrug,scale=fixed,
  options={box range=fixed,min=0,max=20,height=6mm,label=none}]{sample}{箱线图与全部原始样本}
 \gradrow{Atlas,{0,2,2,4,8,16,NA}}
 \gradrow{Birch,{0,0,0,NA}}\gradrow{Cedar,{2,3,4,5,6}}
 \gradrow{Dune,{NA,NA}}
\end{gradtable}\par
rug 的每条短线对应一个有效样本；重复值重叠但计数保留。零是有效值，NA 是缺失；摘要不参与盒列和样本列的统计。
\end{document}
